\section{Further research questions}
\label{sec:future}

The formal classification closes one stated problem, but it also makes
the next tasks more precise. The questions below are not used as
assumptions in any theorem of this article.

\subsection{Short complete certificates for every recurrence pair}

\begin{question}\label{question:certificates}
Given adjacent terms of an even-coefficient recurrence
\eqref{eq:minus-rec} or \eqref{eq:plus-rec}, construct the full
logarithmic correction for $J(a/b)$ with a replayable branch certificate.
Can the number of Hecke and five-term law instances be bounded by a
constant times $\log\max(a,b)$ after consecutive pairs are handled by
Theorem~\ref{thm:neighbor}?
\end{question}

The classification proves existence by rationalized Bloch descent. It
does not automatically make that descent an efficient symbolic
algorithm. A useful implementation should keep the finite log-sine sums
or cyclotomic products in a compressed representation, and state
separately the number of functional-equation steps, the expanded output
size and the bit complexity of algebraic-number operations. An
$O(\log B)$ descent path does not by itself imply an $O(\log B)$
expanded expression.

The neighboring theorem gives a complete base family. The next test
families are the coefficient-$2$ plus sequence
$1,2,5,12,29,\ldots$ and the coefficient-$4$ minus sequence
$1,4,15,56,\ldots$. Their exact classification means a failed numerical
search should now be treated as a defect of the proposed basis or the
certificate construction, not as evidence of formal non-reducibility.

\subsection{Other prime filters}

\begin{question}\label{question:prime-filter}
For a prime $r>2$, classify the formal boundary of
$F(rx)-2F(x)+F(x/r)$ at all positive rational $x$, including the cases
where the $r$-adic valuations of numerator and denominator change the
cyclotomic fields.
\end{question}

At $r=2$, the new layer has an order-two kernel and the four-term witness
in \eqref{eq:four-terms} is decisive. For odd $r$, the analogous new
character space is larger and its sparse group-algebra annihilator must
be analyzed afresh. Simply replacing $2$ by $r$ in
\eqref{eq:mixed-class} has not been justified. A first objective is an
exact conductor-$r$ distribution identity and a precise old-field
intersection theorem.

\subsection{Integral structure and torsion}

\begin{question}\label{question:integral}
What integral lattices underlie the symbol spaces and their dyadic
old/new decompositions? Which denominators and torsion classes are
necessary in explicit five-term certificates?
\end{question}

Rationalization is essential to the clean criterion
\eqref{eq:zero-detection}. It deliberately suppresses the torsion
structure and branch constants that an integral or flattened Bloch-group
implementation must retain. Smith normal forms in a finite
presentation could help locate certificate denominators, but a finite
matrix must not be mistaken for the entire integral pre-Bloch group.

\subsection{Sharper exact counting and uniform depth}

\begin{question}\label{question:counting}
Can the fixed-$R$ expansion \eqref{eq:count-expansion} be made effective
uniformly for $R$ growing with $B$, and can the remaining floor-function
oscillations be isolated by a short exact formula?
\end{question}

The proof gives each fixed fractional-power term and a deliberately
coarse tail. A refined treatment should separate the nearly degenerate
coefficient-$2$ minus family from the geometrically growing families,
and quantify the inversion of the recurrence polynomials uniformly in
their degree. This is a counting problem about the proved formal set,
not a numerical-independence problem.

\subsection{Period evaluation versus formal obstruction}

\begin{question}\label{question:periods}
For which explicitly defined subspaces generated by rational Herglotz
values is numerical period evaluation injective modulo products of
algebraic logarithms and $\pi^2$? Is there an arithmetic statement
strong enough to turn selected nonzero symbols into genuine numerical
non-reduction results?
\end{question}

This question is not answered by the present theorem. The six-extra
classification is a theorem about the rationalized five-term calculus.
Promoting it to an unconditional assertion about all real-number
identities would require an additional theorem of a different kind.
Any conjecture in that direction should first specify the coefficient
field, the permitted logarithmic products and the precise comparison
algebra.

\subsection{A corrected cocycle formulation and colored extensions}

\begin{question}\label{question:cocycle}
Which module and group action, if any, give a valid finite-field cocycle
related to the symbols $\beta_p(a)$ while accommodating the explicit
prime-seven defect? How does that formulation interact with the
Herglotz--Zagier--Novikov parameters of \cite{CK}?
\end{question}

A proposed identity must preserve the nonzero wedge detected in
\eqref{eq:det-witness}, or explicitly state a quotient that kills it.
The rational boundary formula is available independently, so such a
reformulation need not be forced into the proof of the classification.
Colored parameters may require replacing the real positive Gram matrix
by a character-sensitive nondegeneracy argument. This is a distinct
extension, not a consequence of the dyadic proof by notation change.

\subsection{Formal verification targets}

A realistic first formalization would cover the finite group algebra,
normalized norm identities, parity-sensitive boundary expansion and
Vieta descent. The analytic component could then formalize the Rogers
identity on $(0,1)$ and replay the two substitutions in
Section~\ref{sec:identities}. The classical $K$-theory descent and
Dirichlet nonvanishing inputs should remain explicit interfaces until
they are themselves available in the chosen proof assistant.

\section*{Conclusion}
\addcontentsline{toc}{section}{Conclusion}

The all-rational formal $J$ problem is governed by a dyadic old/new
separation, not by a search over increasingly large baskets of constants.
That separation yields a complete modular test, a complete quadratic
recurrence parametrization and a hierarchy of exact height asymptotics.
The analytic Hecke relation then produces a full infinite identity and
a compact exceptional dilogarithm evaluation. The finite-field audit
illustrates why direct symbol calculations, branch control and the
distinction between formal and numerical statements must remain visible
throughout the continuation.
